
\documentclass[11pt]{article}

\usepackage[margin=1in]{geometry}
\usepackage{amsmath,amssymb,amsthm,mathtools,bm}
\usepackage{algorithm}
\usepackage{algpseudocode}
\usepackage{booktabs}
\usepackage{enumitem}
\usepackage{hyperref}

\hypersetup{
  colorlinks=true,
  linkcolor=blue,
  urlcolor=blue,
  citecolor=blue
}

\newcommand{\R}{\mathbb{R}}
\newcommand{\1}{\mathbf{1}}
\newcommand{\Cov}{\operatorname{Cov}}
\newcommand{\diag}{\operatorname{diag}}
\newcommand{\argmin}{\operatorname*{arg\,min}}

\title{Gradient Boosting and Newton-Style Tree Boosting\\
with a Squared Covariance Penalty}
\author{}
\date{}

\begin{document}
\maketitle

\section{Setup}

Consider training data $\{(x_i,y_i)\}_{i=1}^n$, where $y_i$ is the log-price and $f_i := f(x_i)$ is the current prediction. Define the residual vector
\[
e := y - f \in \R^n,
\qquad
e_i = y_i - f_i.
\]
We study the empirical objective
\begin{equation}
\label{eq:main-objective}
\mathcal{L}(f)
=
\frac{1}{2n}\sum_{i=1}^n (y_i-f_i)^2
+
\frac{\rho}{2}\,\widehat{\Cov}(e,y)^2,
\end{equation}
where $\rho \ge 0$ is a regularization parameter and $\widehat{\Cov}(e,y)$ is the empirical covariance between the residuals and the target vector.

We use the convention
\[
\widehat{\Cov}(e,y)
=
\frac{1}{n}\sum_{i=1}^n (e_i-\bar e)(y_i-\bar y),
\]
with
\[
\bar e = \frac{1}{n}\sum_{i=1}^n e_i,
\qquad
\bar y = \frac{1}{n}\sum_{i=1}^n y_i.
\]

Define the centered target vector
\[
c := y - \bar y \1 \in \R^n,
\qquad
c_i = y_i - \bar y.
\]
Since $\sum_{i=1}^n c_i = 0$, we have the exact identity
\begin{equation}
\label{eq:cov-linear-form}
\widehat{\Cov}(e,y)
=
\frac{1}{n} e^\top c
=
\frac{1}{n}(y-f)^\top c.
\end{equation}

Therefore, the penalty term can be rewritten as
\[
\frac{\rho}{2}\left(\frac{1}{n}(y-f)^\top c\right)^2.
\]
This is a \emph{global rank-one quadratic penalty} in the prediction vector $f$, which is the reason the loss becomes nonseparable across observations.

\paragraph{Notation.}
For convenience, define
\begin{equation}
\label{eq:def-z}
z := \widehat{\Cov}(e,y) = \frac{1}{n}(y-f)^\top c.
\end{equation}
Then \eqref{eq:main-objective} becomes
\[
\mathcal{L}(f)
=
\frac{1}{2n}\|y-f\|_2^2 + \frac{\rho}{2} z^2.
\]

\section{Gradient and Hessian with respect to the prediction vector}

Because boosting optimizes over functions through their values on the training set, the relevant derivatives are those with respect to the vector $f=(f_1,\dots,f_n)^\top$.

\subsection{Gradient}

For the quadratic error term,
\[
\nabla_f \left(\frac{1}{2n}\|y-f\|_2^2\right)
=
\frac{1}{n}(f-y).
\]

Now, since
\[
z = \frac{1}{n}(y-f)^\top c,
\]
we get
\[
\nabla_f z = -\frac{1}{n}c.
\]
Hence
\[
\nabla_f \left(\frac{\rho}{2}z^2\right)
=
\rho z \nabla_f z
=
-\frac{\rho}{n}z\,c.
\]

Therefore, the full gradient is
\begin{equation}
\label{eq:gradient}
g := \nabla_f \mathcal{L}(f)
=
\frac{1}{n}(f-y)
-
\frac{\rho}{n} z\,c.
\end{equation}
Componentwise,
\begin{equation}
\label{eq:gradient-components}
g_i
=
\frac{f_i-y_i}{n}
-
\frac{\rho}{n}\,\widehat{\Cov}(e,y)\,(y_i-\bar y).
\end{equation}

Thus, the negative gradient is
\begin{equation}
\label{eq:neg-gradient}
-g_i
=
\frac{y_i-f_i}{n}
+
\frac{\rho}{n}\,\widehat{\Cov}(e,y)\,(y_i-\bar y).
\end{equation}
Up to the common positive factor $1/n$, the pseudo-residual used by first-order boosting is
\begin{equation}
\label{eq:pseudo-residual}
\widetilde r_i
=
e_i + \rho z\,c_i.
\end{equation}

\subsection{Hessian}

Since $z$ is linear in $f$, the Hessian of the penalty term is
\[
\nabla_f^2 \left(\frac{\rho}{2}z^2\right)
=
\rho\,(\nabla_f z)(\nabla_f z)^\top
=
\frac{\rho}{n^2}cc^\top.
\]
Thus the full Hessian is
\begin{equation}
\label{eq:hessian}
H := \nabla_f^2 \mathcal{L}(f)
=
\frac{1}{n}I
+
\frac{\rho}{n^2}cc^\top.
\end{equation}

This is the key structural result:
\begin{itemize}[leftmargin=2em]
\item the gradient is still easy to compute samplewise once the scalar $z$ is known;
\item the Hessian is \emph{not diagonal};
\item instead, it is a diagonal matrix plus a rank-one positive semidefinite correction.
\end{itemize}

\section{Consequences for classical gradient boosting}

\subsection{Functional gradient view}

In standard gradient boosting for squared error, the pseudo-residual is simply the residual $e_i = y_i-f_i$. Here the loss is modified, so the pseudo-residual changes to \eqref{eq:pseudo-residual}:
\[
\widetilde r_i = e_i + \rho z\,c_i.
\]

Therefore, the first-order boosting loop becomes:
\begin{enumerate}[leftmargin=2em]
\item compute current residuals $e_i$;
\item compute the scalar covariance statistic
\[
z = \frac{1}{n}\sum_{i=1}^n e_i c_i;
\]
\item form pseudo-residuals $\widetilde r_i = e_i + \rho z c_i$;
\item fit the next base learner to these pseudo-residuals;
\item compute a line-search step size;
\item update the prediction function.
\end{enumerate}

\subsection{Exact line search for a generic base learner}

Suppose the current base learner has values $h_i = h(x_i)$ on the training set, and consider the update
\[
f^{+} = f + \gamma h.
\]
Then the new residual vector is
\[
e^{+} = y - (f+\gamma h) = e - \gamma h.
\]
By \eqref{eq:cov-linear-form},
\[
\widehat{\Cov}(e^{+},y)
=
\frac{1}{n}(e-\gamma h)^\top c
=
z - \gamma q,
\]
where
\begin{equation}
\label{eq:def-q}
q := \frac{1}{n}c^\top h = \frac{1}{n}\sum_{i=1}^n c_i h_i.
\end{equation}

Therefore, the exact one-dimensional optimization problem is
\begin{equation}
\label{eq:line-search-problem}
\gamma^\star
=
\argmin_{\gamma \in \R}
\left\{
\frac{1}{2n}\|e-\gamma h\|_2^2
+
\frac{\rho}{2}(z-\gamma q)^2
\right\}.
\end{equation}
Expanding and differentiating yields the closed form
\begin{equation}
\label{eq:line-search-closed-form}
\gamma^\star
=
\frac{\frac{1}{n}e^\top h + \rho z q}
{\frac{1}{n}\|h\|_2^2 + \rho q^2}.
\end{equation}

\subsection{Tree base learners: exact leaf-weight update}

Now suppose the weak learner is a regression tree with leaves $I_1,\dots,I_J$ and leaf values $w_1,\dots,w_J$. Then
\[
h(x) = \sum_{j=1}^J w_j \mathbf{1}\{x \in I_j\}.
\]
For a sample $i \in I_j$, the new residual becomes
\[
e_i^{+} = e_i - w_j.
\]
Hence the exact leaf-weight problem for a fixed tree structure is
\begin{equation}
\label{eq:leaf-problem-first-order}
\min_{w \in \R^J}
\left\{
\frac{1}{2n}\sum_{j=1}^J \sum_{i \in I_j} (e_i - w_j)^2
+
\frac{\rho}{2}
\left(
z - \sum_{j=1}^J q_j w_j
\right)^2
\right\},
\end{equation}
where
\[
q_j := \frac{1}{n}\sum_{i \in I_j} c_i.
\]

Define the leaf aggregates
\begin{equation}
\label{eq:leaf-aggregates-first-order}
n_j := |I_j|,
\qquad
b_j := \sum_{i \in I_j} e_i,
\qquad
s_j := \sum_{i \in I_j} c_i,
\qquad
q_j = \frac{s_j}{n}.
\end{equation}
Let
\[
b := (b_1,\dots,b_J)^\top,
\qquad
q := (q_1,\dots,q_J)^\top.
\]
Then the first-order optimality conditions give
\begin{equation}
\label{eq:coupled-leaf-system-first-order}
\left(
\diag\left(\frac{n_1}{n},\dots,\frac{n_J}{n}\right)
+
\rho q q^\top
\right) w
=
\frac{1}{n}b + \rho z q.
\end{equation}

This replaces the usual independent leafwise average formula. The crucial change is that the optimal leaf weights are now \emph{coupled across leaves} through the rank-one matrix $\rho qq^\top$.

\subsection{Efficient closed form via Sherman--Morrison}

Let
\[
D := \diag\left(\frac{n_1}{n},\dots,\frac{n_J}{n}\right),
\qquad
r := \frac{1}{n}b + \rho z q.
\]
Then
\[
w^\star = (D+\rho qq^\top)^{-1}r.
\]
Since this is a diagonal plus rank-one system, Sherman--Morrison gives
\begin{equation}
\label{eq:sherman-morrison-first-order}
(D+\rho qq^\top)^{-1}
=
D^{-1}
-
\frac{\rho D^{-1}qq^\top D^{-1}}
{1+\rho q^\top D^{-1}q},
\end{equation}
and thus
\begin{equation}
\label{eq:leaf-weights-first-order-closed-form}
w^\star
=
D^{-1}r
-
\frac{\rho D^{-1}q \, (q^\top D^{-1}r)}
{1+\rho q^\top D^{-1}q}.
\end{equation}

\section{Newton-style tree boosting (XGBoost / LightGBM style)}

\subsection{Second-order expansion}

In Newton-style tree boosting, one usually builds a second-order approximation of the objective in the additive update $\delta \in \R^n$:
\[
\mathcal{L}(f+\delta)
\approx
\mathcal{L}(f) + g^\top \delta + \frac{1}{2}\delta^\top H \delta.
\]
Here the objective is already quadratic in $f$, so this expansion is in fact \emph{exact}. Using \eqref{eq:gradient} and \eqref{eq:hessian},
\begin{equation}
\label{eq:exact-second-order}
\mathcal{L}(f+\delta)
=
\mathcal{L}(f) + g^\top \delta + \frac{1}{2}\delta^\top
\left(
\frac{1}{n}I + \frac{\rho}{n^2}cc^\top
\right)\delta.
\end{equation}

If the update is constrained to a regression tree with $J$ leaves, write
\[
\delta = Bw,
\]
where $B \in \{0,1\}^{n \times J}$ is the leaf-membership matrix:
\[
B_{ij} = 1
\quad \Longleftrightarrow \quad
x_i \in I_j.
\]
Then the exact tree subproblem becomes
\begin{equation}
\label{eq:newton-tree-subproblem}
\min_{w \in \R^J}
\left\{
g^\top Bw + \frac{1}{2}w^\top(B^\top H B)w + \Omega(w)
\right\},
\end{equation}
with the usual regularizer
\[
\Omega(w) = \gamma J + \frac{\lambda}{2}\|w\|_2^2.
\]

\subsection{Aggregated form}

Define the usual leafwise gradient sums
\begin{equation}
\label{eq:Gj}
G_j := \sum_{i \in I_j} g_i,
\end{equation}
and also
\begin{equation}
\label{eq:sj}
n_j := |I_j|,
\qquad
s_j := \sum_{i \in I_j} c_i.
\end{equation}
Then
\[
B^\top H B
=
\diag\left(\frac{n_1}{n},\dots,\frac{n_J}{n}\right)
+
\frac{\rho}{n^2}ss^\top,
\]
where
\[
s := (s_1,\dots,s_J)^\top.
\]
Therefore, the exact Newton tree subproblem becomes
\begin{equation}
\label{eq:newton-tree-subproblem-aggregated}
\min_{w \in \R^J}
\left\{
G^\top w
+
\frac{1}{2}
w^\top
\left(
\diag\left(\frac{n_1}{n},\dots,\frac{n_J}{n}\right)
+
\frac{\rho}{n^2}ss^\top
+
\lambda I
\right)
w
+
\gamma J
\right\},
\end{equation}
where $G = (G_1,\dots,G_J)^\top$.

Hence the exact optimal leaf weights are
\begin{equation}
\label{eq:newton-leaf-system}
w^\star
=
-
\left(
\diag\left(\frac{n_1}{n},\dots,\frac{n_J}{n}\right)
+
\frac{\rho}{n^2}ss^\top
+
\lambda I
\right)^{-1}
G.
\end{equation}

\paragraph{Comparison with the standard case.}
When $\rho=0$, \eqref{eq:newton-leaf-system} reduces to the familiar independent formula
\[
w_j^\star
=
-\frac{G_j}{\frac{n_j}{n}+\lambda}.
\]
When $\rho>0$, the leaf weights are coupled through the rank-one term $\frac{\rho}{n^2}ss^\top$.

\subsection{Closed form via Sherman--Morrison}

Let
\[
a_j := \frac{n_j}{n}+\lambda,
\qquad
D := \diag(a_1,\dots,a_J),
\qquad
\beta := \frac{\rho}{n^2}.
\]
Then
\[
w^\star = -(D+\beta ss^\top)^{-1}G.
\]
By Sherman--Morrison,
\begin{equation}
\label{eq:sherman-morrison-newton}
(D+\beta ss^\top)^{-1}
=
D^{-1}
-
\frac{\beta D^{-1}ss^\top D^{-1}}
{1+\beta s^\top D^{-1}s}.
\end{equation}
Therefore,
\begin{equation}
\label{eq:newton-leaf-weights-closed-form}
w^\star
=
- D^{-1}G
+
\frac{\beta D^{-1}s \, (s^\top D^{-1}G)}
{1+\beta s^\top D^{-1}s}.
\end{equation}
Componentwise,
\begin{equation}
\label{eq:newton-leaf-weights-componentwise}
w_j^\star
=
-\frac{G_j}{a_j}
+
\frac{\beta s_j}{a_j}
\cdot
\frac{\sum_{k=1}^J \frac{s_k G_k}{a_k}}
{1+\beta \sum_{k=1}^J \frac{s_k^2}{a_k}}.
\end{equation}

\subsection{Exact gain for a fixed tree}

For a fixed tree structure, plugging $w^\star$ into the quadratic objective yields the optimal value
\[
-\frac{1}{2}
G^\top
(D+\beta ss^\top)^{-1}
G
+
\gamma J
\]
up to constants that do not depend on the tree. Therefore, the exact gain is
\begin{equation}
\label{eq:exact-gain}
\operatorname{Gain}(T)
=
\frac{1}{2}
G^\top
(D+\beta ss^\top)^{-1}
G
-
\gamma J.
\end{equation}

This replaces the standard separable gain
\[
\frac{1}{2}\sum_{j=1}^J \frac{G_j^2}{H_j+\lambda} - \gamma J.
\]
The gain is no longer a simple sum of independent leaf contributions, because of the rank-one coupling across leaves.

\section{Why the standard LightGBM/XGBoost custom-objective interface is not exact}

Standard Newton-style boosting libraries assume that the loss can be represented through per-observation first and second derivatives:
\[
\{(g_i,h_i)\}_{i=1}^n,
\]
so that the second-order model is approximated as
\[
\sum_{i=1}^n \left(g_i \delta_i + \frac{1}{2}h_i \delta_i^2\right).
\]
That approximation is exact for separable losses and diagonal Hessians.

In the present setting:
\begin{itemize}[leftmargin=2em]
\item the gradient vector can be computed exactly and returned samplewise;
\item however, the Hessian is \emph{not diagonal}, since
\[
H = \frac{1}{n}I + \frac{\rho}{n^2}cc^\top;
\]
\item therefore, returning only samplewise Hessian entries would ignore the off-diagonal coupling induced by the covariance penalty.
\end{itemize}

As a consequence:
\begin{enumerate}[leftmargin=2em]
\item the usual split-gain formulas are not exact;
\item the usual leaf-weight formula $w_j = -G_j/(H_j+\lambda)$ is not exact;
\item to obtain the correct Newton update, one must modify the tree learner to use the rank-one coupled leaf system in \eqref{eq:newton-leaf-system}.
\end{enumerate}

\section{Algorithmic summaries}

\begin{algorithm}[h!]
\caption{First-order gradient boosting with squared covariance penalty}
\begin{algorithmic}[1]
\Require Training data $\{(x_i,y_i)\}_{i=1}^n$, penalty $\rho \ge 0$, learning rate $\nu \in (0,1]$, number of rounds $M$
\State Compute $c_i \gets y_i - \bar y$, where $\bar y = \frac{1}{n}\sum_{i=1}^n y_i$
\State Initialize $f_0(x) \gets \bar y$
\For{$m = 1,\dots,M$}
    \State Compute current predictions $f_i \gets f_{m-1}(x_i)$ for all $i$
    \State Compute residuals $e_i \gets y_i - f_i$ for all $i$
    \State Compute covariance statistic
    \[
    z \gets \frac{1}{n}\sum_{i=1}^n e_i c_i
    \]
    \State Form pseudo-residuals
    \[
    \widetilde r_i \gets e_i + \rho z c_i
    \qquad \forall i=1,\dots,n
    \]
    \State Fit a base learner structure $T_m$ to $\{(x_i,\widetilde r_i)\}_{i=1}^n$
    \If{$T_m$ is a generic weak learner with output vector $h \in \R^n$}
        \State Compute
        \[
        q \gets \frac{1}{n}c^\top h
        \]
        \State Set exact line-search step
        \[
        \gamma_m \gets
        \frac{\frac{1}{n}e^\top h + \rho z q}
        {\frac{1}{n}\|h\|_2^2 + \rho q^2}
        \]
        \State Update
        \[
        f_m(\cdot) \gets f_{m-1}(\cdot) + \nu \gamma_m h(\cdot)
        \]
    \Else
        \State Let $I_1,\dots,I_J$ be the leaves of $T_m$
        \For{$j=1,\dots,J$}
            \State $n_j \gets |I_j|$
            \State $b_j \gets \sum_{i \in I_j} e_i$
            \State $s_j \gets \sum_{i \in I_j} c_i$
            \State $q_j \gets s_j/n$
        \EndFor
        \State Form
        \[
        D \gets \diag\!\left(\frac{n_1}{n},\dots,\frac{n_J}{n}\right),
        \qquad
        r \gets \frac{1}{n}b + \rho z q
        \]
        \State Solve the coupled leaf system
        \[
        (D+\rho qq^\top)w = r
        \]
        \State Update
        \[
        f_m(\cdot) \gets f_{m-1}(\cdot) + \nu\,T_m(\cdot;w)
        \]
    \EndIf
\EndFor
\State \Return $f_M$
\end{algorithmic}
\end{algorithm}

\begin{algorithm}[h!]
\caption{Exact Newton-style tree boosting with squared covariance penalty}
\begin{algorithmic}[1]
\Require Training data $\{(x_i,y_i)\}_{i=1}^n$, penalty $\rho \ge 0$, leaf regularization $\lambda \ge 0$, complexity penalty $\gamma \ge 0$, learning rate $\nu \in (0,1]$, number of rounds $M$
\State Compute $c_i \gets y_i - \bar y$ for all $i$
\State Initialize prediction function $f_0$
\For{$m = 1,\dots,M$}
    \State Compute current predictions $f_i \gets f_{m-1}(x_i)$
    \State Compute residuals $e_i \gets y_i - f_i$
    \State Compute
    \[
    z \gets \frac{1}{n}\sum_{i=1}^n e_i c_i
    \]
    \State Compute exact gradients
    \[
    g_i \gets \frac{f_i-y_i}{n} - \frac{\rho}{n} z c_i
    \qquad \forall i=1,\dots,n
    \]
    \State Search over candidate tree structures $T$
    \For{each candidate tree $T$ with leaves $I_1,\dots,I_J$}
        \For{$j=1,\dots,J$}
            \State $G_j \gets \sum_{i \in I_j} g_i$
            \State $n_j \gets |I_j|$
            \State $s_j \gets \sum_{i \in I_j} c_i$
        \EndFor
        \State Form
        \[
        D \gets \diag\!\left(\frac{n_1}{n}+\lambda,\dots,\frac{n_J}{n}+\lambda\right),
        \qquad
        \beta \gets \frac{\rho}{n^2},
        \qquad
        A_T \gets D + \beta ss^\top
        \]
        \State Compute exact tree gain
        \[
        \operatorname{Gain}(T)
        \gets
        \frac{1}{2}G^\top A_T^{-1}G - \gamma J
        \]
    \EndFor
    \State Select the best tree structure $T_m$
    \State For the selected tree, form $A_m = D + \beta ss^\top$
    \State Compute exact leaf weights
    \[
    w^\star \gets -A_m^{-1}G
    \]
    \State Update
    \[
    f_m(\cdot) \gets f_{m-1}(\cdot) + \nu\,T_m(\cdot;w^\star)
    \]
\EndFor
\State \Return $f_M$
\end{algorithmic}
\end{algorithm}

\section{Most important points: high-relevance summary}

The most important structural facts are the following.

\begin{enumerate}[leftmargin=2em]
\item \textbf{The covariance penalty is global and nonseparable.}
Because
\[
\widehat{\Cov}(e,y) = \frac{1}{n}(y-f)^\top (y-\bar y\1),
\]
the added regularization term is a rank-one quadratic function of the full prediction vector.

\item \textbf{The pseudo-residual is no longer just the residual.}
For first-order boosting, the correct pseudo-residual is
\[
\widetilde r_i = e_i + \rho\,\widehat{\Cov}(e,y)\,(y_i-\bar y).
\]

\item \textbf{The Hessian is not diagonal.}
The exact Hessian is
\[
H = \frac{1}{n}I + \frac{\rho}{n^2}cc^\top.
\]
This is diagonal plus rank one, not diagonal.

\item \textbf{Tree leaf weights become coupled.}
For a fixed tree, the covariance term prevents independent optimization of each leaf value. The correct leaf values must be obtained from a small linear system of the form
\[
(D + \text{rank-one term})w = \text{right-hand side}.
\]

\item \textbf{The coupling is still computationally cheap.}
Because the extra term is rank one, the exact leaf system can be solved efficiently using Sherman--Morrison, so the correction is mathematically important but computationally manageable.

\item \textbf{Standard XGBoost/LightGBM custom-objective hooks are not exact for this loss.}
They assume per-sample Hessians only. That is insufficient here because the exact Hessian contains off-diagonal terms. Therefore, standard custom objectives can provide the exact gradient but not the exact Newton update unless the tree learner itself is modified.

\item \textbf{There are two clean implementation paths.}
A high-fidelity implementation can be built either as:
\begin{itemize}[leftmargin=2em]
\item an \emph{exact first-order boosting algorithm}, using the modified pseudo-residuals and the exact coupled leaf-weight solve; or
\item an \emph{exact Newton-style tree algorithm}, using the rank-one corrected gain and leaf-weight formulas.
\end{itemize}
\end{enumerate}

\section{Optional implementation remark}

If one wants a practical approximation inside an existing package such as LightGBM or XGBoost without changing the internal split-search code, a reasonable compromise is:
\begin{enumerate}[leftmargin=2em]
\item use the exact gradient \eqref{eq:gradient-components};
\item build the tree structure using the library's standard machinery;
\item once the structure is fixed, replace the library's usual independent leaf-weight formula by the exact coupled solve in \eqref{eq:newton-leaf-system} or \eqref{eq:newton-leaf-weights-closed-form}.
\end{enumerate}
This still does not make the entire procedure fully exact, but it corrects one of the most important consequences of the nonseparable penalty.

\end{document}
